% !TEX root = ../../main.tex

\section{后台管理端技术}

\subsection{Vue 3}

Vue 是一种用于构建用户界面的渐进式 JavaScript 框架，基于标准 HTML、CSS 和 JavaScript 提供声明式、组件化的开发模型\cite{vue-docs}。
Vue 3 在组合式 API、响应式系统和 TypeScript 支持等方面进一步增强，适合构建中后台管理系统和复杂业务页面。

HearBridge 后台管理端采用 Vue 3 技术栈实现，主要用于系统管理员维护手势分类、手势资源、样本数据、训练任务和模型版本。
与移动端强调训练体验不同，后台管理端更关注表格展示、表单编辑、条件查询、状态切换、文件上传和任务触发等管理操作。
Vue 组件化开发方式能够将页面拆分为列表、表单、弹窗、筛选栏、上传组件和状态标签等可复用单元，提高后台页面开发效率。

\subsection{TypeScript}

TypeScript 在 JavaScript 基础上增加了静态类型系统，有助于在开发阶段发现接口字段、组件属性和函数参数方面的错误。
HearBridge 后台管理端涉及较多接口请求和数据表格展示，使用 TypeScript 可以提高前端代码的可维护性，减少字段名写错、数据结构不一致等问题。

在后台管理系统中，接口数据结构相对较多，页面之间存在一定的状态传递和表单校验逻辑。
使用 TypeScript 可以为接口响应、表单数据和页面状态提供更明确的类型约束，从而提高代码可读性和后续维护效率。

\subsection{Vite}

Vite 是现代前端构建工具，具有启动速度快、热更新效率高和构建配置简洁等特点。
后台管理端工程使用 Vite 作为开发和构建工具，并配合类型检查、代码格式化和样式检查等脚本完成日常开发。
对于毕业设计阶段的后台管理系统而言，Vite 能够提高页面开发和调试效率，使开发者更快地完成接口联调和页面迭代。

\subsection{Element Plus}

Element Plus 是基于 Vue 3 的组件库，提供表格、表单、分页、按钮、弹窗、菜单、上传、通知和布局等中后台常用组件\cite{element-plus-docs}。
在后台管理端开发中，使用成熟组件库可以减少基础 UI 组件重复开发工作，使开发重点更多放在业务流程和数据管理上。

HearBridge 后台管理端的核心页面多为典型管理系统页面，例如手势分类列表、手势资源列表、样本管理页面、训练任务页面和模型版本页面。
这些页面通常需要组合使用表格、搜索表单、分页器、编辑弹窗、确认对话框和状态提示。
Element Plus 能够为这些场景提供较完整的组件支持，也有助于保持后台页面风格一致，提升管理员操作效率。
